\clearpage
\phantomsection
\addcontentsline{toc}{chapter}{ABSTRACT}
\chapter*{ABSTRACT}
The project SKILL-GAP \& CURRICULUM DRIFT INTELLIGENCE FOR UNIVERSITIES compares university curricula with industry job postings and the ESCO skills taxonomy. Many of its comparisons depend on a semantic judgement: whether a skill named in a posting and a concept in the taxonomy or a syllabus mean the same thing, one is a narrower kind of the other, they overlap, or they are merely associated. The project's rules resolve most pairs and leave the remainder open. A local language model (Qwen2.5 7B Instruct, run through Ollama) was added to adjudicate those open pairs, together with a deterministic evidence validator and a promotion gate declared before any results were seen. This report evaluates one module of the platform, semantic relationship validation; the other components are described only as context.

The model's answers were compared with stored reviewer responses on a 48-case subset of a frozen 88-case benchmark. Five reviewer records (R01--R05) exist. None qualifies as an independent label set: R02 and R03 consulted an AI assistant before answering, and the answering process of R01, R04 and R05 was not stated. On the 45 cases where R02 and R03 agree (Cohen's $\kappa = 0.92$), the model's raw answer matched the reviewers in 15 cases (33\%). After the evidence validator it matched in 9 cases (20\%). It committed to a definite relationship in 25 cases and was right in 9 of them (36\%). The fixed rules returned \AM{} for every case they were run on and matched none. The model never answered \PA{} or \CH{}. Its main confusions were narrower concepts labelled \SC{} (7 cases), partial overlap labelled \SC{} (5) and broader concepts labelled \RE{} (4). A revised prompt (v4) reduced the over-use of \SC{} but did not raise the validated score, and it was not adopted.

These figures are development estimates on AI-assisted reference answers, not a validated accuracy. The pre-declared gate (semantic-gate-v3) requires two independent human label sets and an accuracy of at least 0.80, so the model remains in an advisory-only Safe Mode: its decisions are recorded and shown but never change a result. The report sets out a staged plan for improving the decision rules, separating development cases from evaluation cases, and collecting independent judgements, with the gate applied unchanged.
\\

\textit{Keywords} : Skill-gap analysis, Curriculum drift, Semantic relationship classification, Curriculum--industry alignment, ESCO, Large language models, Human evaluation, Inter-rater agreement, Selective prediction, Promotion gate
